% ==============================================================================
%  Chapter 1 -- Introduction
%  Purpose: motivate the thesis in ~6-10 pages, readable by a physicist who
%  knows neither ML nor lattice QCD in depth. Written LAST, once all results
%  are in, but the skeleton below fixes the narrative arc now.
% ==============================================================================

\section{Motivation: machine learning for lattice gauge theory}
\label{sec:intro-motivation}

\TODO{Open with the physics problem: non-perturbative QCD observables
(hadron/glueball spectra, topology) are only computable on the lattice, and
lattice computations are dominated by two costs -- ensemble generation and
operator/observable construction. Position ML as a tool for the second:
learning better \emph{operators} and \emph{maps} on gauge configurations,
not replacing the Monte Carlo.}

\TODO{State the central obstruction: gauge symmetry. A generic network fed
raw links must waste capacity learning an exact local symmetry with a group
volume growing like $|\mathcal{G}|^{V}$; a network that is equivariant by
construction gets it for free. Preview the ``inductive-bias gap''
experiment of \cref{sec:reg-inductive-gap} (plaquette-fed CNN
$R^2 \approx 0.99$ vs.\ link-fed CNN $R^2 \approx 0$ on the same target) as
the one-plot motivation.}

\section{From equivariant convolutions to equivariant attention}
\label{sec:intro-attention}

\TODO{Brief history: L-CNN \cite{favoni2022lcnn} showed exact lattice gauge
equivariance via parallel-transported convolutions and bilinear layers;
gauge-covariant ResNets \cite{nagai2021covariant} and the CASK transformer
\cite{cask2025} extended the toolbox. Limitation to stress: convolutional
kernels are \emph{static} after training -- the same stencil everywhere,
on every configuration.}

\TODO{The thesis proposal: GELT, a gauge-equivariant attention encoder.
Attention weights are configuration-dependent, so the network can allocate
its receptive field where the physics is (topological lumps, correlated
regions). Two design departures from L-CNN to state up front: (i) the
value path is matrix-bilinear ($\alpha \, Q^{\dg} \tilde V$), so L-CNN's
loop-growing universality argument transfers; (ii) parallel transport
between sites is averaged over \emph{all} shortest lattice paths in the
$L^1$-ball (a DP recursion, not enumeration).}

\section{Contributions}
\label{sec:intro-contributions}

\TODO{Bullet list, one item per delivered result. Current inventory:}
\begin{itemize}
  \item The GELT architecture (\cref{ch:gelt}): gauge-invariant attention
        scores, matrix-bilinear values, $L^1$-ball shortest-path-averaged
        transport, RoPE geometric prior; equivariance verified by unit
        tests and a numerical invariance check.
  \item A pure-tensor lattice codebase: $\mathbb{Z}_2$/$\SU{2}$ Metropolis
        and $\SU{2}$ heat-bath + overrelaxation samplers (with anisotropy),
        classical glueball toolkit (APE smearing, correlators, GEVP,
        jackknife), and matched-parameter CNN and L-CNN baselines.
  \item \textbf{The headline physics result} (\cref{ch:glueball}): GELT
        trained as a variational $0^{++}$ glueball operator on the Rayleigh
        loss $-C(1)/C(0)$ saturates the transfer-matrix bound, agrees with
        the classical GEVP on the mass ($m\,\at = 0.332 \pm 0.027$), and
        \emph{beats} it on ground-state overlap:
        $A_0 = 0.903 \pm 0.047$ vs.\ $0.837 \pm 0.056$, a correlated
        $\Delta A_0 = +0.078 \pm 0.022$ ($3.6\sigma$) combining the
        original ensemble and an independent replication.
  \item \TODO{The interpretability study (\cref{ch:interpretability}) --
        attention maps as measurements -- once its results exist.}
\end{itemize}

\section{Outline of the thesis}
\label{sec:intro-outline}

\TODO{One sentence per chapter. Ch.~2 lattice gauge theory;
Ch.~3 Monte Carlo methods; Ch.~4 geometric deep learning and prior
equivariant architectures; Ch.~5 the GELT architecture; Ch.~6 validation
and regression experiments; Ch.~7 glueball spectroscopy; Ch.~8
interpretability; Ch.~9 conclusions and outlook.}
